\documentclass{article}
\usepackage[utf8]{inputenc}
\usepackage{graphicx}
\usepackage{amsmath}
\usepackage{booktabs}
\usepackage{float}

\title{Report 1: ECG Heartbeat Categorization}
\author{Nguyen Bich Thu}
\date{\today}

\begin{document}
\maketitle

\begin{abstract}
This report presents a 1D Convolutional Neural Network (1D-CNN) for classifying ECG heartbeats using two public datasets: MIT-BIH Arrhythmia (5 classes) and PTB Diagnostic (2 classes). The model achieves 97.24\% accuracy on MIT-BIH and 91.34\% on PTBDB after 5 training epochs.
\end{abstract}

\section{Introduction}
Electrocardiogram (ECG) signal classification is important for detecting heart abnormalities. This report implements a deep learning approach using 1D-CNN to automatically classify heartbeats from ECG signals.

\section{Dataset Description}
The study uses two datasets from Kaggle:

\subsection{MIT-BIH Arrhythmia Dataset}
\begin{itemize}
    \item \textbf{Task:} Multi-class classification (5 types of heartbeats)
    \item \textbf{Training samples:} 87,554
    \item \textbf{Test samples:} 21,892
    \item \textbf{Class distribution (training set):}
    \begin{itemize}
        \item N (Normal): 72,471 samples (82.8\%)
        \item S: 2,223 samples (2.5\%)
        \item V: 5,788 samples (6.6\%)
        \item F: 641 samples (0.7\%)
        \item Q: 6,431 samples (7.3\%)
    \end{itemize}
\end{itemize}

\subsection{PTB Diagnostic Dataset}
\begin{itemize}
    \item \textbf{Task:} Binary classification (Normal vs Abnormal)
    \item \textbf{Total samples:} Combined from normal and abnormal files
    \item \textbf{Split:} 80\% training, 20\% testing
\end{itemize}

All ECG signals have 187 time-steps, sampled at 125Hz (approximately 1.5 seconds per sample).

\begin{figure}[H]
\centering
\includegraphics[width=0.8\textwidth]{mitbih_samples.png}
\caption{Sample ECG signals from each class in MIT-BIH dataset}
\label{fig:samples}
\end{figure}

\section{Methodology}

\subsection{1D-CNN Architecture}
The model is designed to automatically learn features from raw ECG signals:

\begin{itemize}
    \item \textbf{Input layer:} 187 time-steps × 1 channel
    \item \textbf{Conv1D layer 1:} 64 filters, kernel size 6, ReLU activation
    \item \textbf{BatchNormalization:} Stabilize training
    \item \textbf{MaxPooling1D:} Pool size 3, stride 2
    \item \textbf{Conv1D layer 2:} 64 filters, kernel size 3, ReLU activation
    \item \textbf{BatchNormalization + MaxPooling1D}
    \item \textbf{Flatten + Dense (64 units) + Dropout (0.2)}
    \item \textbf{Output layer:} Softmax (5 units for MIT-BIH, 2 units for PTBDB)
\end{itemize}

\subsection{Training Setup}
\begin{itemize}
    \item \textbf{Optimizer:} Adam
    \item \textbf{Loss function:} Categorical crossentropy
    \item \textbf{Epochs:} 5
    \item \textbf{Batch size:} 64
\end{itemize}

\section{Experiments and Results}

\subsection{MIT-BIH Results}
Training progress over 5 epochs:
\begin{itemize}
    \item Epoch 1: 93.33\% → 96.63\%
    \item Epoch 2: 97.25\% → 96.27\%
    \item Epoch 3: 97.86\% → 97.87\%
    \item Epoch 4: 98.12\% → 97.76\%
    \item Epoch 5: 98.36\% → 97.24\%
\end{itemize}

\begin{table}[H]
\centering
\begin{tabular}{@{}lc@{}}
\toprule
\textbf{Metric} & \textbf{Value} \\
\midrule
Training Accuracy & 98.39\% \\
Validation Accuracy & 97.24\% \\
\bottomrule
\end{tabular}
\caption{MIT-BIH final results}
\label{tab:mitbih}
\end{table}

\subsection{PTBDB Results}
Training progress over 5 epochs:
\begin{itemize}
    \item Epoch 1: 80.62\% → 72.21\%
    \item Epoch 2: 93.40\% → 72.24\%
    \item Epoch 3: 96.78\% → 85.95\%
    \item Epoch 4: 97.10\% → 94.81\%
    \item Epoch 5: 97.91\% → 91.34\%
\end{itemize}

\begin{table}[H]
\centering
\begin{tabular}{@{}lc@{}}
\toprule
\textbf{Metric} & \textbf{Value} \\
\midrule
Training Accuracy & 98.08\% \\
Validation Accuracy & 91.34\% \\
\bottomrule
\end{tabular}
\caption{PTBDB final results}
\label{tab:ptbdb}
\end{table}

\section{Comparison with Original Paper}
According to Kachuee et al. (2018):
\begin{itemize}
    \item MIT-BIH benchmark: ~98\% → Our model: \textbf{97.24\%} (comparable)
    \item PTBDB benchmark: 95.9\% → Our model: \textbf{91.34\%} (slightly lower)
\end{itemize}

The lower result on PTBDB may be due to:
\begin{itemize}
    \item Limited training data
    \item Need for hyperparameter tuning
    \item Different data preprocessing methods
\end{itemize}

\section{Conclusion}
This project successfully implemented a 1D-CNN for ECG heartbeat classification. Key findings:
\begin{itemize}
    \item The model performs well on MIT-BIH dataset (97.24\% accuracy)
    \item Results on PTBDB (91.34\%) show room for improvement
    \item Data imbalance in MIT-BIH (class F only 0.7\%) did not significantly affect results
    \item Future work: experiment with more epochs, different architectures, and data augmentation
\end{itemize}

\section{Citation}
Li, X., Zhang, F., Sun, Z., Li, D., Kong, X., \& Zhang, Y. (2021). Automatic heartbeat classification using S-shaped reconstruction and a squeeze-and-excitation residual network. Computers in Biology and Medicine, 140, 105108. https://doi.org/10.1016/j.compbiomed.2021.105108

\end{document}